Dans le contexte actuel de transformation digitale et de prise de conscience croissante de l'importance de la santé et du bien-être personnel, la capacité à analyser et comprendre nos comportements quotidiens représente un enjeu majeur. Les habitudes de vie telles que le sommeil, l'activité physique, l'alimentation, l'hydratation, le temps d'étude et l'humeur influencent significativement notre qualité de vie et notre performance. Cependant, identifier les anomalies dans ces comportements et comprendre leurs causes reste un défi complexe nécessitant des outils avancés d'analyse de données et d'intelligence artificielle.

Pour répondre à ces enjeux, nous avons conçu et réalisé le \textbf{Personal Behavior Anomaly Detection System (PBADS)}, un système complet de détection d'anomalies comportementales basé sur une architecture microservices. Cette solution intégrée traite les données quotidiennes des utilisateurs depuis leur saisie jusqu'à la génération d'alertes et de recommandations personnalisées, en passant par l'analyse en temps réel à l'aide d'algorithmes d'apprentissage automatique non supervisés.

L'architecture technique de la solution repose sur une approche moderne et robuste. Le système suit une \textbf{architecture microservices} avec 13 services distincts, permettant une séparation claire des responsabilités, une scalabilité élevée et une tolérance aux pannes. Le backend s'appuie sur le framework \textbf{Spring Boot} qui fournit des APIs REST complètes, un système de sécurité avancé basé sur JWT pour l'authentification et l'autorisation, ainsi qu'une intégration avec Kafka pour la communication asynchrone entre services. Le frontend utilise \textbf{Thymeleaf} avec Spring MVC pour le rendu côté serveur, \textbf{Chart.js} pour la visualisation interactive des données, et un système de notifications push en temps réel. Le modèle de données structure de manière cohérente les entités clés du système : \texttt{User} pour la gestion des utilisateurs, \texttt{DailyData} pour les données quotidiennes, \texttt{Alert} pour les alertes générées, et \texttt{Recommendation} pour les recommandations générées par l'IA.

La détection d'anomalies constitue le cœur du système, utilisant des algorithmes d'apprentissage automatique non supervisés (\textbf{Isolation Forest} et \textbf{One-Class SVM}) pour identifier les écarts par rapport aux comportements normaux sans nécessiter de données étiquetées. Le système calcule un score d'anomalie (0-1) et détermine un niveau de sévérité (NORMAL, LOW, MEDIUM, HIGH) pour chaque entrée de données. Lorsqu'une anomalie est détectée, le système génère automatiquement une alerte et sollicite le \textbf{service de recommandations} alimenté par \textbf{Google Gemini AI} pour produire des suggestions personnalisées visant à améliorer les habitudes de l'utilisateur.

La méthodologie de développement a suivi une approche agile avec une planification structurée en sprints, permettant une adaptation rapide aux besoins évolutifs et une livraison incrémentale des fonctionnalités. Cette approche s'est traduite par une priorisation continue des fonctionnalités, des itérations courtes permettant une validation régulière, et des revues régulières pour ajuster les orientations.

Ce rapport technique documente l'ensemble du projet PBADS, couvrant de manière exhaustive le contexte et les objectifs du projet, l'architecture microservices mise en place, la modélisation des données et des processus avec les diagrammes UML, l'implémentation détaillée des microservices, les algorithmes de machine learning utilisés, l'intégration avec Google Gemini AI, et les résultats obtenus. Le document met en évidence les apports concrets de la solution : détection automatique des anomalies comportementales, génération de recommandations personnalisées, notifications en temps réel, et visualisation interactive des données via un tableau de bord moderne avec filtres temporels.

\vspace{1cm}

\textbf{Mots-clés :} Détection d'anomalies, Machine Learning, Microservices, Spring Boot, Kafka, PostgreSQL, Isolation Forest, One-Class SVM, Google Gemini AI, Recommandations personnalisées, Architecture microservices, Thymeleaf, Chart.js, Resilience4j.
